%% Elevark fra mattelab.no – laget automatisk av bygg.py. Rediger gjerne til eget bruk.
%% Kompiler med lualatex.
\documentclass[11pt,a4paper]{article}
\usepackage[norsk]{babel}
\usepackage[a4paper,margin=2.2cm,top=2cm,bottom=2.4cm]{geometry}
\usepackage{fontspec}
\usepackage{unicode-math}
\setmainfont{TeX Gyre Pagella}
\setsansfont{TeX Gyre Heros}[Scale=MatchLowercase]
\setmathfont{TeX Gyre Pagella Math}
\usepackage{xcolor}
\definecolor{blekk}{HTML}{16213A}
\definecolor{dempet}{HTML}{566077}
\definecolor{rute}{HTML}{CFC9B8}
\definecolor{m1}{HTML}{2F4AB0}
\definecolor{m2}{HTML}{99266D}
\usepackage[most]{tcolorbox}
\usepackage{qrcode}
\usepackage{needspace}
\usepackage{fancyhdr}
\usepackage[hidelinks]{hyperref}
\pdfvariable suppressoptionalinfo \numexpr 1+2+4+8+16+32+64+128+256+512\relax
\setlength{\parindent}{0pt}
\setlength{\parskip}{0.4em}
\pagestyle{fancy}
\fancyhf{}
\renewcommand{\headrulewidth}{0pt}
\fancyfoot[L]{\sffamily\footnotesize\color{dempet}mattelab.no · Likningssett løst grafisk · elevark uten fasit}
\fancyfoot[R]{\sffamily\footnotesize\color{dempet}Side \thepage}

\newcommand{\fase}[2]{\tcbox[on line, colframe=#2, colback=white, coltext=#2, boxrule=0.6pt, arc=2pt,
  boxsep=0pt, left=3pt, right=3pt, top=1.5pt, bottom=1.5pt, fontupper=\sffamily\bfseries\footnotesize]{#1}}
\newcommand{\steg}[4]{%
  \par\noindent{\color{rute}\rule{\linewidth}{0.4pt}}\par\nopagebreak
  \noindent\begin{minipage}[t]{2.1cm}\vspace{0pt}\fase{#1}{#2}\end{minipage}%
  \begin{minipage}[t]{\dimexpr\linewidth-2.1cm\relax}\vspace{0pt}#3\end{minipage}\par
  \vspace{#4}}
\newcommand{\oppgave}[2]{\par\Needspace{12\baselineskip}\vspace{0.6em}%
  {\fontsize{14}{17}\selectfont\bfseries{\color{m1}#1}\enspace #2}\par\vspace{0.2em}}

\begin{document}
\begin{minipage}[t]{0.72\linewidth}\vspace{0pt}
{\sffamily\small\color{dempet}Elevark · 1T · mattelab.no}\par\vspace{0.3em}
{\fontsize{24}{28}\selectfont\bfseries Likningssett løst grafisk}\par\vspace{1em}
{\sffamily\small Navn: \rule{4.6cm}{0.4pt}\quad Klasse: \rule{1.5cm}{0.4pt}\quad Dato: \rule{1.9cm}{0.4pt}}
\end{minipage}\hfill
\begin{minipage}[t]{0.22\linewidth}\vspace{0pt}\raggedleft
\qrcode[height=2.3cm]{https://mattelab.no/utforskere/likningssett.html}\par\vspace{0.2em}
{\sffamily\scriptsize\color{dempet}Utforskeren}
\end{minipage}

\vspace{1em}
{\sffamily\small\color{dempet}Utforskeren ligger på \url{https://mattelab.no/utforskere/likningssett.html}}\par
Hver oppgave har tre steg. \textbf{Gjett} først, uten å bruke utforskeren. \textbf{Test} så i utforskeren. \textbf{Forklar} til slutt med egne ord hvorfor det blir slik.

\oppgave{1}{Les av skjæringspunktet}
\steg{Gjett}{m1}{Se på likningssettet $2x + y = 5$ og $x - y = 1$. Prøv deg fram med små hele tall. Hvilke $x$ og $y$ passer i begge likningene?}{2.2cm}
\steg{Test}{m2}{Sett opp likningssettet. Les av skjæringspunktet, og se på kontrollen i panelet.}{2.6cm}
\steg{Forklar}{blekk}{Hvorfor er skjæringspunktet løsningen av likningssettet? Løs likningssettet med addisjonsmetoden.}{3.4cm}

\oppgave{2}{Parallelle linjer}
\steg{Gjett}{m1}{Se på $x + 2y = 4$ og $x + 2y = 1$. Hvor mange løsninger tror du likningssettet har? Hva skjer hvis den andre likningen er $x + 2y = 4$? Eller $2x + 4y = 8$?}{2.2cm}
\steg{Test}{m2}{Sett opp likningssettet. Dra $c_2$ sakte opp til $4$. Prøv deretter $2x + 4y = 8$ som likning 2.}{2.6cm}
\steg{Forklar}{blekk}{Prøv å løse $x + 2y = 4$ og $x + 2y = 1$ ved regning. Hva skjer? Hvorfor har $x + 2y = 4$ og $2x + 4y = 8$ uendelig mange løsninger?}{3.4cm}

\oppgave{3}{Når grafen ikke er nøyaktig nok}
\steg{Gjett}{m1}{Linjene $3x + 2y = 6$ og $x - 2y = 1$ skjærer hverandre et sted mellom $x = 1$ og $x = 2$. Hvor tror du skjæringspunktet er?}{2.2cm}
\steg{Test}{m2}{Sett opp likningssettet. Zoom inn rundt skjæringspunktet. Hvor nøyaktig klarer du å lese det av?}{2.6cm}
\steg{Forklar}{blekk}{Løs likningssettet med regning. Skriv svaret som brøk og som desimaltall.}{3.4cm}

\oppgave{4}{Kinobilletter}
\steg{Gjett}{m1}{Tre voksne og to barn betaler $590$ kr for kino. En voksen og fire barn betaler $480$ kr. La $x$ være prisen for en voksen og $y$ prisen for et barn. Sett opp to likninger. Hva tror du billettene koster?}{2.2cm}
\steg{Test}{m2}{Sett opp likningssettet. Trykk «Vis hele grafen» hvis du ikke ser skjæringspunktet.}{2.6cm}
\steg{Forklar}{blekk}{Løs likningssettet med innsettingsmetoden. Kontroller at svaret passer i teksten.}{3.4cm}

\oppgave{5}{Når finnes det ingen løsning?}
\steg{Gjett}{m1}{Se på likningssettet $kx + y = 3$ og $2x + y = 1$. For hvilken verdi av $k$ tror du det ikke finnes noen løsning?}{2.2cm}
\steg{Test}{m2}{Sett opp likningssettet med $k = 0$. Dra $a_1$ (som er $k$) sakte oppover. Følg med på skjæringspunktet og på tallet $a_1b_2 - a_2b_1$.}{2.6cm}
\steg{Forklar}{blekk}{Bruk stigningstallene til å forklare svaret. Hva skjer med skjæringspunktet når $k$ nærmer seg denne verdien?}{3.4cm}

\end{document}
